% GENERATED FILE. Edit canonical lessons, then run npm run generate:books.
% canonical-source-hash: fnv1a64:81aa74fe00cb12cf
% canonical-lessons: FR-C400-peripherique, FR-C400-peripherie, FR-C400-echangeur, FR-C400-bretelle, FR-C400-bifurquer

\chapter{Autour de la ville en voiture}
\label{ch:fr-autour-de-la-ville-en-voiture}
\hlchaptermodality{400}

\begin{chapteropening}
\textbf{By the end of this chapter:} \emph{I can talk about the big roads round a town: the ring road, the outskirts, the interchange, the slip road and turning off.}

Complete the last lesson of chapter 400: I can talk about the big roads round a town: the ring road, the outskirts, the interchange, the slip road and turning off.
\end{chapteropening}

\section[le périphérique]{le périphérique — the ring road (round a city)}

\label{lesson:FR-C400-peripherique}



\begin{quote}
Before the new one: say the French for a team, then the French for a birthday.
\end{quote}

\subsection*{You'll want to know: le périphérique}
\textbf{le périphérique} — \textquotedblleft{}the ring road (round a city)\textquotedblright{}.

\begin{quote}
Le soir, il y a beaucoup de voitures sur le périphérique.
\end{quote}

\begin{quote}
\emph{In the evening there are a lot of cars on the ring road.}
\end{quote}

In everyday speech people often say \emph{le périph}.

\begin{cousinweb}
\emph{périphérique} is built on \emph{périphérie} 'outer edge', a learned word that is usually traced through Late Latin \emph{peripheria} to Greek \emph{periphéreia} 'circumference', from \emph{peri} 'around' and \emph{phérein} 'to carry'. EtymDB has no reliable source line for it.
\end{cousinweb}

\begin{grammarlens}[title={What you've built}]
One more word for telling, explaining and describing at B1.
\end{grammarlens}

\subsection*{Your turn}
\begin{itemize}
\raggedright
  \item \emph{Say it:} \emph{le périphérique}
  \item \emph{Say it:} \emph{le périphérique}, once more
  \item \emph{Say it:} a sentence of your own with \emph{le périphérique}
  \item \emph{Recall:} say the French for a team, then the French for a birthday, then say \emph{le périphérique} again
\end{itemize}

\begin{tcolorbox}[breakable,colback=teal!4,colframe=teal!35!black,title={Before you move on}]
What does le périphérique mean? (\textquotedblleft{}The ring road (round a city)\textquotedblright{}.) Say it once more.
\end{tcolorbox}
\section[la périphérie]{la périphérie — the outskirts, the edge (of a town)}

\label{lesson:FR-C400-peripherie}



\begin{quote}
Before the new one: say the French for a birthday, then the French for the ring road.
\end{quote}

\subsection*{You'll want to know: la périphérie}
\textbf{la périphérie} — \textquotedblleft{}the outskirts, the edge (of a town)\textquotedblright{}.

\begin{quote}
Mes cousins habitent à la périphérie de la ville.
\end{quote}

\begin{quote}
\emph{My cousins live on the outskirts of the town.}
\end{quote}

\emph{En périphérie} — on the edge of town, away from the centre.

\begin{cousinweb}
\emph{périphérie} is usually traced to Late Latin \emph{peripheria}, a learned borrowing from Greek \emph{periphéreia} 'circumference', built from \emph{peri} 'around' and \emph{phérein} 'to carry'. EtymDB does not list the word, so the line is hedged.
\end{cousinweb}

\begin{grammarlens}[title={What you've built}]
One more word for telling, explaining and describing at B1.
\end{grammarlens}

\subsection*{Your turn}
\begin{itemize}
\raggedright
  \item \emph{Say it:} \emph{la périphérie}
  \item \emph{Say it:} \emph{la périphérie}, once more
  \item \emph{Say it:} a sentence of your own with \emph{la périphérie}
  \item \emph{Recall:} say the French for a birthday, then the French for the ring road, then say \emph{la périphérie} again
\end{itemize}

\begin{tcolorbox}[breakable,colback=teal!4,colframe=teal!35!black,title={Before you move on}]
What does la périphérie mean? (\textquotedblleft{}The outskirts, the edge (of a town)\textquotedblright{}.) Say it once more.
\end{tcolorbox}
\section[l'échangeur]{l'échangeur — the interchange (between motorways)}

\label{lesson:FR-C400-echangeur}



\begin{quote}
Before the new one: say the French for the ring road, then the French for the outskirts.
\end{quote}

\subsection*{You'll want to know: l'échangeur}
\textbf{l'échangeur} — \textquotedblleft{}the interchange (between motorways)\textquotedblright{}.

\begin{quote}
Après l'échangeur, nous avons pris l'autoroute du sud.
\end{quote}

\begin{quote}
\emph{After the interchange we took the southbound motorway.}
\end{quote}

A masculine noun: \emph{un échangeur}.

\begin{cousinweb}
\emph{échangeur} is built on \emph{échanger} 'to exchange' with the ending \emph{-eur}. \emph{Échanger} comes from Old French \emph{eschangier}, which is inherited from Vulgar Latin \emph{excambiāre} (EtymDB): the place where cars exchange one road for another.
\end{cousinweb}

\begin{grammarlens}[title={What you've built}]
One more word for telling, explaining and describing at B1.
\end{grammarlens}

\subsection*{Your turn}
\begin{itemize}
\raggedright
  \item \emph{Say it:} \emph{l'échangeur}
  \item \emph{Say it:} \emph{l'échangeur}, once more
  \item \emph{Say it:} a sentence of your own with \emph{l'échangeur}
  \item \emph{Recall:} say the French for the ring road, then the French for the outskirts, then say \emph{l'échangeur} again
\end{itemize}

\begin{tcolorbox}[breakable,colback=teal!4,colframe=teal!35!black,title={Before you move on}]
What does l'échangeur mean? (\textquotedblleft{}The interchange (between motorways)\textquotedblright{}.) Say it once more.
\end{tcolorbox}
\section[la bretelle]{la bretelle — the slip road; the strap}

\label{lesson:FR-C400-bretelle}



\begin{quote}
Before the new one: say the French for the outskirts, then the French for the interchange.
\end{quote}

\subsection*{You'll want to know: la bretelle}
\textbf{la bretelle} — \textquotedblleft{}the slip road; the strap\textquotedblright{}.

\begin{quote}
Le camion a pris la bretelle pour quitter l'autoroute.
\end{quote}

\begin{quote}
\emph{The lorry took the slip road to leave the motorway.}
\end{quote}

\emph{Les bretelles} are also braces that hold up trousers.

\begin{cousinweb}
\emph{bretelle} was borrowed from Old High German \emph{brittil} 'rein, strap' (EtymDB). A strap ties two things together, and so does a short road between two big ones.
\end{cousinweb}

\begin{grammarlens}[title={What you've built}]
One more word for telling, explaining and describing at B1.
\end{grammarlens}

\subsection*{Your turn}
\begin{itemize}
\raggedright
  \item \emph{Say it:} \emph{la bretelle}
  \item \emph{Say it:} \emph{la bretelle}, once more
  \item \emph{Say it:} a sentence of your own with \emph{la bretelle}
  \item \emph{Recall:} say the French for the outskirts, then the French for the interchange, then say \emph{la bretelle} again
\end{itemize}

\begin{tcolorbox}[breakable,colback=teal!4,colframe=teal!35!black,title={Before you move on}]
What does la bretelle mean? (\textquotedblleft{}The slip road; the strap\textquotedblright{}.) Say it once more.
\end{tcolorbox}
\section[bifurquer]{bifurquer — to turn off, to fork}

\label{lesson:FR-C400-bifurquer}



\begin{quote}
Before the new one: say the French for the interchange, then the French for the slip road.
\end{quote}

\subsection*{You'll want to know: bifurquer}
\textbf{bifurquer} — \textquotedblleft{}to turn off, to fork\textquotedblright{}.

\begin{quote}
Au village, la route bifurque vers la montagne.
\end{quote}

\begin{quote}
\emph{At the village the road forks towards the mountains.}
\end{quote}

\emph{La bifurcation} — the fork in the road.

\begin{cousinweb}
\emph{bifurquer} is a learned formation on Latin \emph{bifurcus} 'with two prongs', from \emph{bi-} 'two' and \emph{furca} 'fork' (Lewis \& Short).
\end{cousinweb}

\begin{grammarlens}[title={What you've built}]
One more word for telling, explaining and describing at B1.
\end{grammarlens}

\subsection*{Your turn}
\begin{itemize}
\raggedright
  \item \emph{Say it:} \emph{bifurquer}
  \item \emph{Say it:} \emph{bifurquer}, once more
  \item \emph{Say it:} a sentence of your own with \emph{bifurquer}
  \item \emph{Recall:} say the French for the interchange, then the French for the slip road, then say \emph{bifurquer} again
\end{itemize}

\begin{tcolorbox}[breakable,colback=teal!4,colframe=teal!35!black,title={Before you move on}]
What does bifurquer mean? (\textquotedblleft{}To turn off, to fork\textquotedblright{}.) Say it once more.
\end{tcolorbox}
